% Complete technical proofs for the recovery and information theorems.
% Input this fragment after the manuscript's appendix command.

\section{Why the recovery threshold is sharp}\label{rec:threshold-proof}

The construction has three stages. Estimate the common error rate;
denoise the table using its finite-prime approximations; then turn
approximate frequencies into exact rationals. The last stage can be
very slow. The theorem asserts consistency and gives no useful practical
sample bound for the particular implementation below.

\subsection{Calibrating the noise without knowing its sign}

Write $J$ for the all-ones matrix and $I_N$ for the identity. All data
matrices have zero diagonal. Set
\[
 \widehat q=\frac1{N(N-1)}\sum_{i\ne j}Y_{ij},\qquad
 \widehat a=\operatorname{clip}_{[-1,1]}
 \left(\frac{2\widehat q-1}{2\kappa-1}\right).
\]
Here $\operatorname{clip}_{[u,v]}$ means projection onto $[u,v]$.
If $q_0$ is the clean empirical edge density, conditioning on $E$ gives
\[
 \widehat q=\tfrac12+a_N(q_0-\tfrac12)+O_{\Prob}(N^{-1}).
\]
Changing one integer changes $q_0$ by at most $2/N$, whence
$q_0-\E q_0=O_{\Prob}(N^{-1/2})$. With $X=N^2$, Möbius inversion gives
\begin{equation}\label{rec:finite-kappa}
 \E q_0=\frac1{X^2}\sum_{d\le X}\mu(d)\lfloor X/d\rfloor^2
 =\kappa+O(\log X/X).
\end{equation}
Replacing squared floors costs $O(X^{-1}\sum_{d\le X}d^{-1})$;
the omitted tail of $\sum_d\mu(d)/d^2$ costs $O(X^{-1})$.
Thus, for nonzero $a_N$, clipping cannot increase the absolute error and
\begin{equation}\label{rec:calibration-error}
 \frac{\widehat a}{a_N}-1
 =O_{\Prob}\left(N^{-1/2}+\frac1{N|a_N|}+\frac{\log N}{N^2}\right).
\end{equation}
When $T_N\to\infty$, this tends to zero, the sign is eventually correct
with probability tending to one, and $N\widehat a^2/T_N\to1$ in
probability. In the independent-prime model the finite-range bias in
\eqref{rec:finite-kappa} is absent.

\subsection{The matrix estimate}

For a matrix $U$, denote its operator and Frobenius norms by
$\|U\|_{\mathrm{op}}$ and $\|U\|_F$. We record the exact deterministic
estimate needed, so no uniformity condition is hidden in an external
denoising theorem.

\begin{lemma}[Spectral truncation]\label{rec:svt}
Suppose $B=A+Z$ are symmetric and $\|Z\|_{\mathrm{op}}\le t/2$.
Retain the eigenvalues of $B$ with absolute value at least $t$, with
their signs, to obtain $\widehat A$. For every matrix $R$ of rank at
most $r$,
\[
 \|\widehat A-A\|_F^2
 \le C\bigl(rt^2+\|A-R\|_F^2\bigr)
\]
for an absolute constant $C$.
\end{lemma}
\begin{proof}
Let $s_i$ be the decreasing singular values of $A$ and let $r_0$ count
those at least $t/2$. Weyl's inequality bounds the rank of $\widehat A$
by $r_0$. If $A_0$ is the rank-$r_0$ truncation of $A$, then
\[
 \|\widehat A-A_0\|_{\mathrm{op}}
 \le\|\widehat A-B\|_{\mathrm{op}}+\|Z\|_{\mathrm{op}}
       +\|A-A_0\|_{\mathrm{op}}\le2t.
\]
The difference has rank at most $2r_0$, giving
$\|\widehat A-A\|_F^2\le16r_0t^2+2\sum_{i>r_0}s_i^2$.
For every $r$,
\[
 r_0t^2\le rt^2+4\sum_{i>r}s_i^2,
 \qquad
 \sum_{i>r_0}s_i^2\le\sum_{i>r}s_i^2+rt^2/4.
\]
The best-rank approximation inequality
$\sum_{i>r}s_i^2\le\|A-R\|_F^2$ completes the proof.
\end{proof}

Apply this to
\[
 B=Y-\tfrac12(J-I_N),\qquad
 A=a_NM,\qquad M=E-\tfrac12(J-I_N),
\]
with threshold $t=10\sqrt N$. Conditional on the integers, the entries
of $Z=B-A$ above the diagonal are independent, centered, and supported
on intervals of length one. A $1/4$-net of the unit sphere has at most
$9^N$ points. For a fixed unit vector $v$, Hoeffding's inequality gives
$\Prob(|v^TZv|>u\mid X_1,\ldots,X_N)\le2e^{-u^2}$, because the sum
of squared interval lengths in $2\sum_{i<j}v_iv_jZ_{ij}$ is at most
two. The net inequality therefore yields
\begin{equation}\label{rec:op-concentration}
 \Prob(\|Z\|_{\mathrm{op}}>4\sqrt N\mid X_1,\ldots,X_N)
 \le2\exp\bigl(-(4-\log9)N\bigr).
\end{equation}

The arithmetic input is a low-rank approximation. Let $F^{(m)}_{ij}$,
including the diagonal, indicate that $i,j$ share none of the first
$m$ primes. There are only $2^m$ truncated vertex profiles, so
\[
 \operatorname{rank}(F^{(m)}-J/2)\le2^m+1.
\]
Off the diagonal this differs from $M$ only if a larger prime is shared.
The probability of sharing $p$ is at most $p^{-2}$, also for ordinary
integers, and $p_m\ge m+1$. Consequently
\begin{equation}\label{rec:rank-tail}
 \E\frac{\|M-(F^{(m)}-J/2)\|_F^2}{N^2}
 \le\frac1m+\frac1{4N}.
\end{equation}
The last term accounts for the diagonal, whose approximation error is
exactly $1/2$ in absolute value.

For the proof choose $m=\lfloor\tfrac12\log_2T_N\rfloor$.
This choice is not used by the algorithm. Lemma~\ref{rec:svt},
\eqref{rec:op-concentration}, \eqref{rec:rank-tail}, and Markov's
inequality give
\[
 \frac{\|\widehat A-a_NM\|_F^2}{a_N^2N^2}
 =O_{\Prob}\left(T_N^{-1/2}+\frac1{\log T_N}\right).
\]
For $i\ne j$, define
$\widehat E_{ij}=\operatorname{clip}_{[0,1]}(1/2+\widehat A_{ij}/\widehat a)$,
and set the diagonal to zero. At $\widehat a=0$ use an arbitrary fixed
table. Equation~\eqref{rec:calibration-error} and clipping show
\begin{equation}\label{rec:frobenius}
 Q_N:=\frac1{N^2}\|\widehat E-E\|_F^2
 =O_{\Prob}((\log T_N)^{-1}).
\end{equation}

\subsection{Approximate frequencies become exact primes}

The algorithm commutes with vertex permutations. Thus
\eqref{rec:frobenius} controls every fixed row in probability, even
though it is initially a statement about the whole matrix. To see the
rate carefully, let $G_C=\{Q_N\le C/\log T_N\}$. This event is
permutation invariant. On it, the expected squared error of any fixed
row, averaged across its columns and multiplied by $\mathbf1_{G_C}$,
is $\E[Q_N\mathbf1_{G_C}]\le C/\log T_N$. Choose $C$ to make $G_C$
likely, then use Markov's inequality on $G_C$. Hence for fixed $i$,
\begin{equation}\label{rec:row-error}
 \frac1N\sum_k(\widehat E_{ik}-E_{ik})^2
 =O_{\Prob}((\log T_N)^{-1}).
\end{equation}

Form $\widehat d_i=N^{-1}\sum_k\widehat E_{ik}$ and
$\widehat c_{ij}=N^{-1}\sum_k\widehat E_{ik}\widehat E_{jk}$.
Cauchy--Schwarz bounds their denoising errors by
$O_{\Prob}((\log T_N)^{-1/2})$. For the second assertion use
\[
 |uv-u'v'|\le |u-u'|+|v-v'|\quad(0\le u,v,u',v'\le1).
\]
For ordinary integers put $D(n)=\varphi(n)/n$ and
$L_{ij}=\operatorname{lcm}(X_i,X_j)$. Inclusion--exclusion gives
\begin{equation}\label{rec:inclusion-exclusion}
 \left|\Prob_{U\le X}((U,n)=1)-D(n)\right|
 \le\frac{2^{\omega(n)}}X.
\end{equation}
For every fixed $\delta>0$, $2^{\omega(n)}\le C_\delta n^\delta$:
in the product $2^{\omega(n)}/n^\delta$ only finitely many primes can
give a factor above one. Take $\delta=1/4$ and note $L_{ij}\le X^2$,
$X=N^2$. The bias in \eqref{rec:inclusion-exclusion} is uniformly
$O(N^{-1})$. Conditional on $X_i,X_j$, the other vertices are
independent, so their empirical clean degree and codegree errors are
$O_{\Prob}(N^{-1/2})$. Their targets are $D(X_i),D(X_j),D(L_{ij})$.

Division by the last target is controlled uniformly, since
\begin{equation}\label{rec:codegree-tight}
 \E[-\log D(L_{ij})]
 \le2\sum_p\frac{-\log(1-1/p)}p<\infty.
\end{equation}
Thus $D(L_{ij})^{-1}=O_{\Prob}(1)$. The corresponding infinite-model
products satisfy the same estimate. With clipping and a default at a
zero estimated denominator, we have in both models
\begin{equation}\label{rec:estimated-ratio}
 \widehat r_{ij}:=
 \operatorname{clip}_{[0,1]}\left(\frac{\widehat d_i\widehat d_j}
                                    {\widehat c_{ij}}\right)
 =\frac{\varphi(R_{ij})}{R_{ij}}
   +O_{\Prob}((\log T_N)^{-1/2}).
\end{equation}

Set $\widehat T=N\widehat a^2$ and
\[
 H_N=\max\left\{1,\left\lfloor
       [\log(\max\{e,\widehat T\})]^{1/16}\right\rfloor\right\}.
\]
Choose the nearest reduced rational to $\widehat r_{ij}$ in $[0,1]$
with denominator at most $H_N$, using a fixed rule for ties. If it is
a finite totient ratio, apply Lemma~\ref{rec:totient-decoder}; otherwise
return a fixed default. This is a finite test: reject zero, reject a
largest denominator prime of exponent exceeding one, remove that prime
as prescribed, and reject if the ratio leaves $(0,1]$ or the largest
denominator prime fails to decrease. Accept only on reaching one.

The true denominator is tight. For ordinary integers it is at most
$\gcd(X_i,X_j)$ and, for $H\ge1$,
\begin{equation}\label{rec:gcd-tail}
 \Prob(\gcd(X_i,X_j)>H)
 \le\sum_{d>H}\Prob(d\mid X_i,\ d\mid X_j)
 \le\sum_{d>H}d^{-2}=O(H^{-1}).
\end{equation}
In the independent-prime model $R_{ij}$ is a finite random integer,
which gives the same tightness requirement without a rate.
Calibration implies $H_N\asymp(\log T_N)^{1/16}$ in probability.
Distinct candidate rationals differ by at least $H_N^{-2}$, while
\[
 H_N^2\left|\widehat r_{ij}-\frac{\varphi(R_{ij})}{R_{ij}}\right|
 =O_{\Prob}((\log T_N)^{-3/8})\longrightarrow0.
\]
The nearest rational is therefore correct with probability tending to
one. Permutation equivariance makes this the expected fraction of
incorrect pair outputs; Markov's inequality proves the fraction
assertion in Theorem~\ref{rec:threshold}.

\subsection{A single hidden parity bit proves necessity}

Suppose along a subsequence that $T_N\le C$, so eventually $|a_N|\le1/2$.
Giving a decoder additional information can only help it. We give almost
everything away and retain one fair bit that changes whether a pair
shares the prime 2.

With $X=N^2$, the disjoint pairs $\{u,2u\}$, $u\le X/2$ odd, cover
a fraction $1/2+o(1)$ of the possible values of $X_i$. On this event
reveal the odd base $u$, leaving a fair choice between $u$ and $2u$.
Reveal every other integer and restrict also to $X_j$ being even. This
joint event has probability $1/4+o(1)$. The correct $R_{ij}$ contains
2 precisely in the second case.

Only the at most $N-1$ edges incident to $i$ can distinguish the cases.
Each differing observation exchanges Bernoulli parameters
$(1-a_N)/2$ and $(1+a_N)/2$, whose relative entropy is
\[
 a_N\log\frac{1+a_N}{1-a_N}\le4a_N^2\qquad(|a_N|\le1/2).
\]
Conditional independence bounds the total relative entropy by $4T_N$.
For two laws $P,Q$ with equal priors, the testing error is at least
$e^{-D(P\Vert Q)}/4$. Indeed, their Hellinger affinity $\rho$ satisfies
$\rho\ge e^{-D(P\Vert Q)/2}$ by Jensen, while
$\operatorname{TV}(P,Q)\le\sqrt{1-\rho^2}$ by Cauchy--Schwarz.
The error $(1-\operatorname{TV})/2$ is consequently at least $\rho^2/4$.
Every pair estimator thus has error at least
\begin{equation}\label{rec:testing-lower}
 (1/16+o(1))e^{-4T_N}.
\end{equation}
In the independent-prime model, hide only $B_{i,2}$ and require $B_{j,2}=1$;
the same argument gives $e^{-4T_N}/8$.
The bounds are uniform in the choice of pair. Their positive lower bound
on a bounded-$T_N$ subsequence also bounds the expected pair-error
fraction. A random variable in $[0,1]$ tending to zero in probability
must have expectation tending to zero. Necessity follows.

\section{Counting the information that survives}\label{rec:information-proof}

We prove Theorem~\ref{rec:information} directly for ordinary integers.
This avoids approximating a growing array of divisibility bits by
independent prime bits. The upper bound asks how much information small
primes can carry and how much the remaining primes can transmit through
the noise. The lower bound temporarily reveals a small set of reference
integers and uses them to identify many prime features of the others.

\subsection{The prime sums and the conditioning principle}

Write $h(t)=-t\log t-(1-t)\log(1-t)$, with its endpoint values defined
by continuity, and put
\[
 S(y)=\sum_{p\le y}h(1/p),\qquad \rho_y=\sum_{p>y}p^{-2}.
\]
The classical prime estimates give, for $y\ge2$,
\begin{equation}\label{rec:prime-estimates}
 \begin{split}
 S(y)&=\log y+\log\log y+O(1),\\
 \sum_{p\le y}\frac{\log p}{p}&=\log y+O(1),\qquad
 \sum_{p\le y}\log p=O(y),\qquad
 \rho_y\ll\frac1{y\log y}.
 \end{split}
\end{equation}
The first follows by expanding
$h(1/p)=(\log p)/p+1/p+O(p^{-2})$ and using Mertens' sums.
The last follows by partial summation from $\pi(y)\ll y/\log y$.
These are classical estimates, with stronger explicit versions in
Rosser and Schoenfeld~\cite{RecRosserSchoenfeld1962}; see also
Mertens~\cite{RecMertens1874} for the prime sums.

We shall use one elementary information inequality repeatedly. If $A$
and $Z$ are any data determined by the hidden integers, then the channel
noise is independent of them given $G=G_N$, so
\begin{equation}\label{rec:oracle-information}
 \begin{split}
 I(G;Y)&=I(G,A;Y)
       =I(A;Y)+I(G;Y\mid A)\\
       &\ge I(G;Y\mid A)\ge I(Z;Y\mid A).
 \end{split}
\end{equation}
The last step is conditional data processing. Thus reference values may
be revealed \emph{in a proof of a lower bound} without claiming they
were given to the actual decoder. This particular inequality depends
on the conditional independence of the observation channel.

\subsection{Upper bound: truncate the prime features}

Reveal the array $B$ of prime-divisibility bits for all vertices and
all $p\le y$. The bits of one integer are dependent, but entropy
subadditivity still gives
\begin{equation}\label{rec:bit-entropy}
 H(B)\le N\sum_{p\le y}h\left(\frac{\lfloor X/p\rfloor}{X}\right)
 \le NS(y),\qquad X=N^2.
\end{equation}
Here $h$ is increasing on $[0,1/2]$ and
$\lfloor X/p\rfloor/X\le1/p\le1/2$.

Given $B$, let $F_{ij}$ be the answer that uses only primes at most
$y$. If $F_{ij}=0$, the clean answer is determined. If $F_{ij}=1$,
the clean answer differs only if a prime above $y$ divides both hidden
integers. The average probability of this event is at most $\rho_y$.

First assume $|a_N|\le1/2$. The divergence between the two output laws
for an edge is at most $4a_N^2$, as in the testing proof. For a binary
input that differs from a fixed baseline with probability $r$, its
mutual information through this channel is at most $4a_N^2r$.
To verify this, compare the two output distributions with the baseline
output law: the average divergence is $rD$, and subtracting the
nonnegative divergence of the mixture from the baseline gives mutual
information. Conditional entropy subadditivity across edges, together
with independent noise given $G$, now yields
\begin{equation}\label{rec:info-upper-cutoff}
 I(G;Y)\le H(B)+I(G;Y\mid B)
 \le NS(y)+4a_N^2\binom N2\rho_y.
\end{equation}
Only the averaged residual probability is used here; independence of
prime divisibility for ordinary integers is unnecessary.

If $T_N\ge2$, choose $y=T_N$. Equations
\eqref{rec:prime-estimates} and \eqref{rec:info-upper-cutoff} give
\[
 I(G;Y)\le N\log T_N+N\log\log T_N+O(N).
\]
For bounded $T_N$, choosing $y=2$ instead gives $I(G;Y)=O(N)$,
uniformly on each bounded range of $T_N$.

If $|a_N|>1/2$, then $T_N>N/4$. We use a clean entropy bound. Reveal
prime bits through $N$ and then record which edges differ from the
truncated graph. Entropy subadditivity and concavity of $h$ give
\begin{equation}\label{rec:clean-upper}
 H(G)\le NS(N)+\binom N2h(\rho_N)
 \le N\log N+N\log\log N+O(N),
\end{equation}
because $h(\rho_N)=O(N^{-1})$. Since $I(G;Y)\le H(G)$ and
$\log(1+T_N)=\log N+O(1)$ in this regime, these estimates prove the
upper half of \eqref{rec:information-law} for every noise rate.

\subsection{Lower bound: reference integers reveal many small primes}

Put $L=\log N$, $m=\lfloor N/L\rfloor$, and
\begin{equation}\label{rec:reference-cutoff}
 y=T_N/L^{12}.
\end{equation}
If $y<2$, then $\log(1+T_N)=O(\log\log N)$, and the required lower
bound follows from $I(G;Y)\ge0$. Suppose henceforth $y\ge2$.
Reveal the actual values of the first $m$ sampled integers, calling
this reference data $A$. The remaining $N-m$ vertices are targets.

For a target value $n$, call it good when $2^{\omega(n)}\le L^3$.
The divisor identity $2^{\omega(n)}=\sum_{d\mid n}\mu(d)^2$ gives
\[
 \frac1X\sum_{n\le X}2^{\omega(n)}
 \le\sum_{d\le X}\frac1d\le1+\log X.
\]
Thus a uniformly sampled target is bad with probability $O(L^{-2})$.
For each prime $p\le y$, use all reference integers divisible by $p$.
Let their number be $L_p$. It is binomial with mean at least $m/(2p)$.
A Chernoff bound gives
\begin{equation}\label{rec:reference-count}
 \Prob(L_p<m/(4p))\le e^{-cm/p}
\end{equation}
for an absolute $c>0$ and all sufficiently large $N$.

If $p\mid n$, every clean answer from this target to these references
is zero. If $p\nmid n$ and $n$ is good, inclusion--exclusion gives
\begin{align}
 \#\{b\le X:p\mid b,\ (b,n)=1\}
 &=\frac Xp D(n)+O(2^{\omega(n)})\notag\\
 &\ge\frac X{pL^3}-L^3
 \ge\frac X{2pL^3}.\label{rec:reference-separation}
\end{align}
We used $D(n)\ge2^{-\omega(n)}$ and
$X/p\ge N L^{12}$, since $p\le T_N/L^{12}\le N/L^{12}$.
Consequently a random reference conditioned to be divisible by $p$
is coprime to this good $n$ with probability at least $1/(2L^3)$.

Conditional on $n$ and on which reference indices are divisible by
$p$, their values are independent and uniform over the multiples of
$p$ in $[1,X]$. Their noisy answers are therefore independent. When
$p\mid n$ their mean is $\varepsilon_N$; otherwise it is
$\varepsilon_N+a_N\theta$, with $\theta\ge1/(2L^3)$ on good targets.
Use the threshold $\varepsilon_N+a_N/(4L^3)$, reversing the comparison
when $a_N<0$. The parameter is fixed in the probability law and may
be used by this auxiliary estimator for the information lower bound.
The reconstruction algorithm in Theorem~\ref{rec:threshold} does not
require it.

On the event in \eqref{rec:reference-count}, Hoeffding's inequality
bounds the bit error on a good target by
\begin{equation}\label{rec:reference-testing}
 2\exp\left(-c\frac{a_N^2L_p}{L^6}\right)
 \le2\exp\left(-c\frac{T_N}{pL^7}\right)
 \le2e^{-cL^5}.
\end{equation}
The count failure probability is smaller after modifying $c$.
The independence calculation conditions on the reference divisibility
indicators, \emph{not} on their full revealed values. We average over
those values to obtain the displayed error bound. The resulting
estimator is nevertheless a function of $(Y,A)$, as required in
\eqref{rec:oracle-information}.

Let $D_i^*=\prod_{p\le y,\ p\mid X_i}p$ be a target's truncated
squarefree mask. Union-bound over the at most $N$ candidate primes.
The probability of failing to recover the whole mask is
\begin{equation}\label{rec:mask-error}
 e_N=O(L^{-2})+O(Ne^{-cL^5}).
\end{equation}
There are at most $X$ possible masks, since each is an integer divisor
of some number at most $X$. Replace an impossible estimate by 1. The
usual error-flag proof of Fano's inequality gives
\begin{equation}\label{rec:fano-mask}
 H(D_i^*\mid Y,A)\le h(e_N)+e_N\log X=O(L^{-1}).
\end{equation}
For completeness, condition additionally on the indicator that the
estimated mask is wrong. On a correct estimate there is no uncertainty;
on an incorrect one there are at most $X$ choices. The entropy of the
indicator is at most $h(e_N)$.

There is a simple lower bound for the mask's own entropy. For every
possible value $d$, $\Prob(D_i^*=d)\le\Prob(d\mid X_i)\le1/d$.
Therefore, using \eqref{rec:prime-estimates},
\begin{align}
 H(D_i^*)&\ge\E\log D_i^*
 =\sum_{p\le y}\log p\,\frac{\lfloor X/p\rfloor}{X}\notag\\
 &\ge\sum_{p\le y}\frac{\log p}{p}
       -\frac1X\sum_{p\le y}\log p
 \ge\log y-O(1).\label{rec:mask-entropy}
\end{align}
The target masks are mutually independent and independent of $A$.
Apply \eqref{rec:oracle-information} to their collection and subtract
their conditional entropy using \eqref{rec:fano-mask}. We obtain
\begin{align}
 I(G;Y)&\ge (N-m)(\log y-O(1))-O(N/L)\notag\\
 &\ge N\log T_N-12N\log\log N-O(N).\label{rec:info-lower}
\end{align}
Here $m\log y=O(N)$ because $y\le N$. Replacing $T_N$ by $1+T_N$
costs $O(N)$ when $T_N\ge1$; for smaller $T_N$ the lower bound was
already trivial. This proves \eqref{rec:information-law}. Finally
take $\varepsilon_N=0$, so $Y=G$ and $T_N=N$, to obtain
\eqref{rec:graph-entropy}. This derivation of the clean entropy uses
no earlier graph-entropy theorem.

The proof shows why the leading scale is logarithmic in $T_N$.
The entropy in prime bits through $y$ is roughly $\log y$ per vertex,
and reference testing accesses these bits up to $T_N$ with only
logarithmic losses. The coefficient comes from arithmetic prime sums.
General entropy behavior of random-free graph models is studied by
Hatami and Norine~\cite{RecHatamiNorine2012}; graph mutual information
and signal-dependent information laws have broader precedents in
Skeja--Olhede~\cite{RecSkejaOlhede2024} and
Deshpande--Abbe--Montanari~\cite{RecDeshpandeAbbeMontanari2015}.
The assertion here is the specific finite, labelled coprimality law.


\section{Proof of reconstruction without digit addresses}\label{rec:message-proof}

\begin{proof}[Proof of Theorem~\ref{rec:message}]
By Corollary~\ref{rec:profiles}, for each fixed $j$ we can form an
estimated lane $\widehat L_j$ differing from
$L_j=\{i:\lpf(X_i)=p_j\}$ at $o_{\Prob}(N)$ vertices. The exceptional
value $X_i=1$ belongs to no lane. Ordinary residue counting and the
law of large numbers give
\begin{equation}\label{rec:lane-size}
 \frac{|L_j|}{N}\longrightarrow
 \delta_j=\frac1{p_j}\prod_{q<p_j}(1-1/q)>0.
\end{equation}
Use
\begin{equation}\label{rec:lane-estimator}
 \widehat h_j=\frac1{|\widehat L_j|}
              \sum_{i\in\widehat L_j}\frac{1-2Z_i}{b_N},
\end{equation}
with a fixed default when the denominator vanishes.

Conditional on the integers and edge noise, the reconstructed lane
is fixed and the digit errors remain independent. Each summand in
\eqref{rec:lane-estimator} has conditional mean $1-2x_g(X_i)$,
whose absolute value is at most one, and conditional variance at
most $b_N^{-2}$. Equations \eqref{rec:lane-size} and $Nb_N^2\to\infty$
make the centered error $o_{\Prob}(1)$. Replacing the estimated lane
by the true lane in the \emph{conditional mean} costs only
$o_{\Prob}(1)$. No division of this replacement error by $|b_N|$
is needed.

The average Liouville sign in each fixed lane tends to zero. To spell
out the classical input, write $Q_j=\prod_{q<p_j}q$ and
$L(t)=\sum_{n\le t}\lambda(n)=o(t)$, the prime number theorem in
its Liouville form. Inclusion--exclusion and complete multiplicativity
give
\[
 \sum_{m\le t,\ (m,Q_j)=1}\lambda(m)
 =\sum_{d\mid Q_j}\mu(d)\lambda(d)L(t/d)
 =\sum_{d\mid Q_j}L(t/d)=o(t).
\]
The least-prime lane consists of $p_jm$ with $(m,Q_j)=1$; multiplication
by $p_j$ reverses the sign. This proves that its finite-range mean
of $1-2x_g(n)$ tends to $1-g_j$. The independently sampled lane
average has the same limit by bounded-variance concentration and
\eqref{rec:lane-size}. Thus $\widehat h_j\to1-g_j$ in probability.
Rounding at $1/2$ recovers the bit. A union bound recovers any fixed
finite prefix of mask bits.

Define, with arbitrary binary defaults where needed,
\[
 \widehat C_N=\sum_{n=2}^N
 \frac{1-\lambda(n)}2\widehat g_{\ell(n),N}\,2^{-n}.
\]
For every fixed $M$, the first $M$ digits are correct with probability
tending to one; the remaining tail is at most $2^{-M}$. First taking
$N\to\infty$ and then $M\to\infty$ proves convergence to $C_g$.

For necessity of the digit condition, compare two fixed masks differing
only at some $j\ge2$. Supply an oracle with the exact integers, every
other mask bit, and both noise probabilities. The two hypotheses
differ at at most $N$ digit observations. On a subsequence where
$Nb_N^2$ is bounded, their total conditional relative entropy is
at most $4Nb_N^2$ eventually. The testing bound used in
\eqref{rec:testing-lower} leaves positive error under at least one
mask. Thus one procedure cannot be consistent for both.

For the edge condition, use the all-one mask. The earlier pair
$u,2u$ has opposite Liouville signs, so it would be distinguished by
a clean digit label. Instead use the disjoint pairs
$\{u,4u\}$ with $u$ odd and $u\le N^2/4$. Their union has probability
$1/4+o(1)$, and the two values have exactly the same Liouville sign.
Reveal the base and every other hidden integer and require the partner
to be even. This event has probability $1/8+o(1)$ and leaves a fair
bit that changes membership of 2 in $R_{ij}$ but changes no digit
observation law. The incident edges again have total divergence at
most $4Na_N^2$. On a bounded-$Na_N^2$ subsequence the fixed-pair error
is bounded away from zero, uniformly over pairs, which also excludes
a vanishing error fraction. This proves the joint necessity statement.
\end{proof}
